\documentclass[12pt,a4paper]{article}
\usepackage[T1]{fontenc}\usepackage[utf8]{inputenc}\usepackage{times}
\usepackage[a4paper,top=2.5cm,bottom=2.5cm,left=2.5cm,right=2.5cm]{geometry}
\usepackage{graphicx}
\graphicspath{{figures/}}
\usepackage{booktabs}\usepackage{array}\usepackage{enumitem}\usepackage{tabularx}
\usepackage{caption}\usepackage{fancyhdr}\usepackage{parskip}\usepackage{float}
\usepackage{setspace}\onehalfspacing
\usepackage{xcolor}
\usepackage{xurl}\usepackage[hidelinks]{hyperref}
\definecolor{navy}{HTML}{1B4F72}
\newcommand{\F}[1]{{\small\texttt{#1}}}
\captionsetup{font=small,labelfont=bf,labelsep=colon,
              justification=justified,singlelinecheck=false}
\setlist[itemize]{leftmargin=1.3em,itemsep=4pt,topsep=4pt}
\setlist[enumerate]{leftmargin=1.6em,itemsep=6pt,topsep=4pt}
\pagestyle{fancy}\fancyhf{}
\fancyhead[L]{\footnotesize KOUAME Koffi Fid\`ele}
\fancyhead[C]{\footnotesize HEALTHCARE NO-SHOWS}
\fancyhead[R]{\footnotesize Data Analysis Internship}
\fancyfoot[C]{\thepage}\renewcommand{\headrulewidth}{0.4pt}
\setlength{\headheight}{15pt}
\setlength{\emergencystretch}{2em}

% A workbook capture: framed with a hairline so that it sits cleanly on the page
\newcommand{\capture}[3][\linewidth]{%
\begin{figure}[H]\centering
\setlength{\fboxsep}{0pt}\setlength{\fboxrule}{0.4pt}%
\fbox{\includegraphics[width=\dimexpr#1-0.8pt\relax]{#2}}
\caption{#3}
\end{figure}}
% Two charts of the dashboard side by side
\newcommand{\pair}[3]{%
\begin{figure}[H]\centering
\setlength{\fboxsep}{0pt}\setlength{\fboxrule}{0.4pt}%
\fbox{\includegraphics[width=0.485\linewidth]{#1}}\hfill
\fbox{\includegraphics[width=0.485\linewidth]{#2}}
\caption{#3}
\end{figure}}
% A clickable entry of the cover
\newcommand{\go}[2]{\hyperref[#1]{#2}}

\begin{document}

\begin{titlepage}\centering
\singlespacing
\vspace*{3cm}
{\large\scshape Data Analysis Internship}\\[2.6cm]
\rule{\linewidth}{0.5pt}\\[0.8cm]
{\LARGE\bfseries Healthcare No-Shows}\\[0.6cm]
{\large Why Patients Miss Their Appointments:\\[2pt] Data Cleaning, KPIs and an Interactive Excel Dashboard}\\[0.8cm]
\rule{\linewidth}{0.5pt}\\[1.2cm]
{\normalsize \go{sec:data}{Data} $\bullet$ \go{sec:quality}{Data Quality} $\bullet$
\go{sec:cleaning}{Cleaning} $\bullet$ \go{sec:kpis}{KPIs}\\[4pt]
\go{sec:results}{Results} $\bullet$ \go{sec:insights}{Insights} $\bullet$
\go{sec:dashboard}{Excel Dashboard}}\\[3.2cm]
{\large\itshape Prepared by}\\[0.5cm]
{\Large\bfseries KOUAME Koffi Fid\`ele}\\[0.6cm]
{\large \today}
\vfill\end{titlepage}

\singlespacing
\tableofcontents\newpage
\onehalfspacing

% =====================================================================
\section{Executive summary}\label{sec:summary}

The task is to turn the raw file \F{healthcare\_noshows.csv} into a professional,
interactive Excel dashboard: clean and prepare the data, define the relevant KPIs, build
the charts and extract the insights that explain why patients miss their appointments.
The whole work is done in Excel, in one workbook: \F{Healthcare\_NoShows\_Dashboard.xlsx}.

The file holds \textbf{106{,}987 appointments} held between 29 April and 8 June 2016.
After cleaning, \textbf{106{,}977 appointments} of \textbf{60{,}268 patients} remain, and
\textbf{21{,}672} of them were missed: \textbf{one appointment in five (20.3\%)}.

\textbf{Waiting is the main driver.} An appointment made for the same day is missed in
\textbf{4.7\%} of cases; an appointment booked one day or more ahead in \textbf{28.5\%},
six times more.

\textbf{The SMS reminder works, although the raw figures say the opposite.} Patients who
received an SMS miss more often (27.7\% against 16.7\%), because the reminder is only
sent for appointments booked several days ahead. Compared at the same lead time, the
SMS lowers the no-show rate in every band; the Mantel-Haenszel odds ratio computed in
the workbook is \textbf{0.78}.

\textbf{Past behaviour predicts the next visit.} After a previous no-show the rate is
\textbf{29.4\%}, against 17.1\% after an attended visit. Teenagers (26.4\%) and welfare
recipients (+3.9 points) also miss more; the weekday makes almost no difference.

The file supports a simple targeting rule: remind every patient booked one day or more
ahead, and call the patients who already missed an appointment.

% =====================================================================
\section{Data and method}\label{sec:data}

\subsection{The data}

The file supplied with the task has 15 columns and one row per appointment: the patient
and appointment identifiers (\F{PatientId}, \F{AppointmentID}), the gender and age, the
neighbourhood of the health unit, the booking and appointment dates (\F{ScheduledDay},
\F{AppointmentDay}) and their difference in days (\F{Date.diff}), five health and social
flags (\F{Scholarship}, \F{Hipertension}, \F{Diabetes}, \F{Alcoholism}, \F{Handcap}),
the reminder flag \F{SMS\_received} and the outcome \F{Showed\_up}.

\subsection{Method}

The work follows the brief step by step, with Excel tools only.

\begin{enumerate}
\item \textbf{Audit:} completeness, duplicates, ranges and chronology of the dates
(Excel table, COUNTIFS, COUNTBLANK).
\item \textbf{Cleaning:} invalid rows removed, columns renamed, codes recoded (recoding
tables, PROPER, SUBSTITUTE).
\item \textbf{Preparation:} nine analysis fields added by formula (IF, CHOOSE, WEEKDAY).
\item \textbf{KPIs:} nine indicators, each answering one question (AVERAGE, SUM,
AVERAGEIFS).
\item \textbf{Analysis:} rates by segment and chi-square tests of the differences
(COUNTIFS, SUMIFS, CHISQ.DIST.RT).
\item \textbf{Dashboard and insights:} two interactive dashboards (pivot tables, pivot
charts, slicers, drop-down lists) and seven findings, each with an action and an owner.
\end{enumerate}

Every number of this report is read from the workbook, where it is a live formula on
the clean table \F{tbl\allowbreak Appointments}.

% =====================================================================
\section{Data quality}\label{sec:quality}

\begin{table}[H]\centering\singlespacing
\caption{Audit of the raw file (106{,}987 rows)}
\small
\begin{tabularx}{\linewidth}{@{}>{\raggedright\arraybackslash}Xr>{\raggedright\arraybackslash}p{5.8cm}@{}}
\toprule
\textbf{Check} & \textbf{Rows} & \textbf{Decision} \\ \midrule
Empty cells & 0 & None found \\
Duplicate AppointmentID & 0 & None found \\
Date.diff different from the two dates & 0 & None found \\
Appointment dated before its booking & \textbf{5} & \textbf{Removed}: impossible chronology \\
Age of 115 years & \textbf{5} & \textbf{Removed}: implausible age \\
PatientId with decimals & 5 & Kept, stored as text \\
Second booking of a patient on the same day & 8{,}532 & Kept: own identifier and outcome \\
Appointments on a Saturday & 39 & Kept: rare but valid \\
Misspelled column names & 3 & Renamed \\
Codes TRUE/FALSE, F/M, upper-case names & 106{,}987 & Recoded \\
\bottomrule
\end{tabularx}
\end{table}

The file is complete and internally consistent: no cell is empty, no appointment is
duplicated, and \F{Date.diff} always equals the difference between the two dates. Two
defects make a row unusable. Five appointments are dated before their booking, which is
impossible; five others belong to two patients recorded at 115 years, an implausible
age. The other findings are valid situations (several bookings on one day, rare Saturday
clinics) or presentation issues.

% =====================================================================
\section{Cleaning and preparation}\label{sec:cleaning}

\subsection{Cleaning}

The ten invalid rows are removed and listed in full on the \emph{Data Cleaning} sheet,
so that the removal can be checked. Three columns are renamed (Hipertension to
Hypertension, Handcap to Handicap, Date.diff to LeadTimeDays). TRUE/FALSE become Yes/No,
F/M become Female/Male, and \F{Showed\_up} becomes \F{Status} (Attended, No-show). The
neighbourhood names, written in capitals, are converted with the formula
\F{PROPER}, the particles da, de, do and das being set back in lower case with
\F{SUBSTITUTE}; the 81 names and their formula are shown on the same sheet.

\subsection{Preparation}

Nine fields are added to the Excel table \F{tblAppointments}, each by a formula copied
down the column, so that the table stays consistent if a value is corrected.

\begin{table}[H]\centering\singlespacing
\caption{Fields added by Excel formulas}
\small
\begin{tabularx}{\linewidth}{@{}l>{\raggedright\arraybackslash}X@{}}
\toprule
\textbf{Field} & \textbf{Definition} \\ \midrule
NoShow & 1 for a no-show, 0 otherwise: its average is directly the no-show rate \\
LeadTimeDays & AppointmentDate minus ScheduledDate \\
LeadTimeBand & Same day, 1 to 3, 4 to 7, 8 to 14, 15 to 30, 31 to 60, 61+ days (nested IF) \\
AgeGroup & 00-11, 12-17, 18-24, 25-34, 35-44, 45-54, 55-64, 65-74, 75+ (nested IF) \\
Weekday, WeekStart & Day name (CHOOSE and WEEKDAY) and Monday of the appointment week \\
SMSFlag & 1 if an SMS was received: its average is the SMS coverage \\
PriorNoShows & Number of earlier no-shows of the same patient \\
PatientHistory & First appointment, Attended before or Previous no-show \\
\bottomrule
\end{tabularx}
\end{table}

\F{PriorNoShows} compares each row with the row above: the table is sorted by patient,
then date, then appointment number, and the formula adds the no-shows of the patient's
earlier rows. \F{PatientHistory} follows from it.

\subsection{Checks}

Live formulas on the clean table confirm the result: 106{,}987 raw rows minus 10 removed
rows equal the 106{,}977 rows of the table; no cell is empty; no lead time is negative;
the oldest patient is 102 years old; and 60{,}268 rows are a first appointment, which is
the number of patients.

\capture{data_cleaning.png}{Data Cleaning sheet: audit of the raw file, the ten rows
removed, the recoding table and, on the right, the live checks on the clean table.}

% =====================================================================
\section{KPIs}\label{sec:kpis}

Each KPI answers one operational question. The table gives its value on the clean table
and the Excel formula that computes it on the \emph{KPI Definitions} sheet.

\begin{table}[H]\centering\singlespacing
\caption{Key performance indicators}
\small
\begin{tabularx}{\linewidth}{@{}l>{\raggedright\arraybackslash}Xr@{}}
\toprule
\textbf{KPI} & \textbf{Question answered} & \textbf{Value} \\ \midrule
Appointments & How much demand is there? & 106{,}977 \\
Missed appointments & How many slots are lost? & 21{,}672 \\
No-show rate & How often is a slot wasted? & \textbf{20.3\%} \\
Average lead time & How far ahead are patients booked? & 10.2 days \\
Same-day share & How much activity is booked on the day? & 34.7\% \\
SMS coverage & How many patients are reminded? & 32.3\% \\
No-show rate, booked ahead & What is the risk once a patient waits? & 28.5\% \\
Repeat no-show rate & Does past behaviour predict the next visit? & 29.4\% \\
Attendance rate & How reliable is the schedule? & 79.7\% \\
\bottomrule
\end{tabularx}
\end{table}

The last two rate KPIs name their own base: the booked-ahead rate is computed on the
69{,}824 appointments booked one day or more ahead, the repeat rate on the 15{,}563
appointments of patients who had already missed one.

% =====================================================================
\section{Results}\label{sec:results}

The charts of this section are those of the \emph{Dashboard} sheet, with every filter set
to All.

\subsection{Lead time}

\capture[0.8\linewidth]{chart_lead_time.png}{No-show rate by lead time. Same-day
appointments are 34.7\% of the volume but only 8.0\% of the no-shows.}

As soon as the patient has to wait, the rate jumps from 4.7\% to 22.9\% for 1 to 3 days
and climbs to 34.4\% for 31 to 60 days. The slight fall beyond 60 days concerns only
2{,}044 appointments.

\subsection{The SMS reminder}

\capture{sms_analysis.png}{SMS Analysis sheet: no-show rate with and without SMS in each
lead-time band (same day excluded, since no SMS is sent), Mantel-Haenszel odds ratio and
slots recoverable.}

Overall, the SMS group misses 10.9 points more. This is a Simpson's paradox. No SMS is
sent for same-day appointments, and only 6.2\% of the 1 to 3 day appointments receive
one, against about 60\% beyond three days: the SMS group is made of long-lead
appointments, which carry the highest risk. Compared like with like, the SMS lowers the
rate in every band, by 1.6 points (1 to 3 days) to 7.8 points (61+ days). The
Mantel-Haenszel odds ratio, which pools the six bands, is \textbf{0.78}: at equal lead
time the reminder cuts the odds of a no-show by about a fifth.

\subsection{Who misses appointments}

\pair{chart_age.png}{chart_history.png}{No-show rate by age group (left) and by patient
history (right). Teenagers of 12 to 17 years miss 26.4\% of their appointments, patients
of 65 to 74 years 15.1\%. A patient who already missed an appointment misses the next one
in 29.4\% of cases, against 17.1\% after an attended visit.}

\subsection{Health and social factors}

\begin{table}[H]\centering\singlespacing
\caption{No-show rate with and without each factor (chi-square test, 1 degree of freedom)}
\small
\begin{tabular}{@{}lrrrrr@{}}
\toprule
\textbf{Factor} & \textbf{Group size} & \textbf{Rate} & \textbf{Others} & \textbf{Difference} & \textbf{$p$-value} \\ \midrule
Scholarship (welfare) & 10{,}809 & 23.8\% & 19.9\% & $+3.9$ & $<0.001$ \\
Hypertension & 21{,}800 & 17.3\% & 21.0\% & $-3.7$ & $<0.001$ \\
Diabetes & 7{,}943 & 18.0\% & 20.4\% & $-2.4$ & $<0.001$ \\
Handicap & 2{,}234 & 18.0\% & 20.3\% & $-2.3$ & 0.007 \\
Alcoholism & 3{,}360 & 20.1\% & 20.3\% & $-0.1$ & 0.872 \\
Gender (female) & 70{,}110 & 20.4\% & 20.1\% & $+0.3$ & 0.293 \\
\bottomrule
\end{tabular}
\end{table}

Welfare recipients miss 3.9 points more. Gender and alcoholism make no difference.
Hypertension and diabetes seem protective, but hypertensive patients are mostly older,
and older patients miss less. Compared at equal age, the gap disappears: from 45 years,
where most hypertensive patients are, their rate is equal to or higher than that of the
others.

\capture{risk_factors.png}{Risk Factors sheet: rates with and without each factor with
the chi-square test, hypertension at equal age, and the neighbourhood indicators.}

\subsection{When and where}

\pair{chart_weekday.png}{chart_week.png}{No-show rate by weekday (left) and by week of
the appointment (right). From Monday to Friday the rate moves between 19.5\% and 21.3\%
(the 39 Saturday appointments are too few to count); over the six weeks it stays between
18\% and 22\%.}

\capture[0.8\linewidth]{chart_neighbourhoods.png}{The ten neighbourhoods with the most
missed appointments. Together they account for 37.2\% of all no-shows.}

Place matters more than time. Among the 53 neighbourhoods with at least 500
appointments, the no-show rate ranges from 15.7\% to 29.1\%, a spread of thirteen
points, against less than two points between weekdays.

% =====================================================================
\section{Insights and recommendations}\label{sec:insights}

\subsection{Lead time is the primary risk}

4.7\% of same-day appointments are missed, against 28.5\% of the appointments booked one
day or more ahead.

\textbf{Recommendation.} Keep a share of same-day and short-notice slots, and apply a
moderate overbooking to the slots booked more than a week ahead, where nearly one
appointment in three is missed. \emph{Owner: scheduling.}

\subsection{The reminder works but misses its target}

At equal lead time the SMS lowers the odds of a no-show by about a fifth. Yet 35{,}240
appointments booked one day or more ahead received no SMS, and 29.4\% of them were missed.

\textbf{Recommendation.} Send the reminder to every appointment booked one day or more
ahead, starting with the 1 to 3 day bookings, of which only 6.2\% are reminded today. If
these appointments had the no-show rate observed with an SMS in their band, about
\textbf{1{,}400 slots} would have been saved over the six weeks. \emph{Owner: patient
communication.}

\subsection{Past no-shows repeat}

A patient who already missed an appointment misses the next one in 29.4\% of cases,
against 17.1\% after an attended visit.

\textbf{Recommendation.} Flag these patients at booking and call them the day before the
appointment. \emph{Owner: front desk.}

\subsection{Young patients and welfare recipients need a different contact}

Teenagers of 12 to 17 years miss 26.4\% of their appointments and young adults of 18 to
24 years 25.0\%. Welfare recipients miss 3.9 points more than other patients.

\textbf{Recommendation.} Address the reminders for minors to their parents and offer
slots outside school hours. Review access barriers (transport, working hours) with the
social services. \emph{Owner: scheduling and social services.}

\subsection{Chronic patients are not safer}

The lower rate of hypertensive patients is an age effect, not a behaviour.

\textbf{Recommendation.} Do not reduce the reminders sent to chronic patients.
\emph{Owner: care coordination.}

\subsection{Place, not the weekday, is the lever}

The weekday moves the rate by less than two points, the neighbourhood by thirteen, and
ten neighbourhoods hold 37\% of the no-shows.

\textbf{Recommendation.} Keep the weekly schedule unchanged and focus outreach on the ten
neighbourhoods that lose the most slots. \emph{Owner: public health.}

\capture[0.8\linewidth]{insights.png}{Insights sheet: each insight with its evidence, a live value
computed on the clean table, the recommended action and its owner.}

% =====================================================================
\section{The Excel dashboard}\label{sec:dashboard}

\F{Healthcare\_NoShows\_Dashboard.xlsx} opens on a \emph{Home} sheet whose contents link
to every other sheet.

\begin{table}[H]\centering\singlespacing
\caption{Sheets of the workbook}
\small
\begin{tabularx}{\linewidth}{@{}l>{\raggedright\arraybackslash}X>{\raggedright\arraybackslash}p{5.6cm}@{}}
\toprule
\textbf{Sheet} & \textbf{Content} & \textbf{Excel tools} \\ \midrule
Home & Contents with a link to each sheet & Hyperlinks \\
Dashboard & 6 drop-down filters, 6 KPI cards, 8 charts & Data validation, COUNTIFS, SUMIFS \\
Pivot Dashboard & 6 slicers, 6 KPI cards, 8 pivot charts & Slicers, pivot charts, GETPIVOTDATA \\
Insights & Seven insights with evidence, live value, action and owner & AVERAGEIFS \\
Risk Factors & Factors with chi-square tests, hypertension at equal age, neighbourhoods & COUNTIFS, CHISQ.DIST.RT, MINIFS, data bars \\
SMS Analysis & SMS at equal lead time, Mantel-Haenszel odds ratio & COUNTIFS, conditional formatting \\
Dashboard Data & Formula tables behind the Dashboard, dynamic top 10 & INDEX, MATCH, LARGE \\
Pivot Tables & The pivot tables behind the Pivot Dashboard & Pivot tables, top 10 filter \\
Data Cleaning & Audit, rows removed, recoding, formulas, checks & PROPER, SUBSTITUTE, COUNTBLANK \\
KPI Definitions & Question, definition, formula and value of each KPI & Structured references \\
Clean Data & 106{,}977 appointments, 23 columns & Excel table \\
\bottomrule
\end{tabularx}
\end{table}

The workbook offers two ways to filter. On the \emph{Dashboard}, six drop-down lists
(gender, age group, lead time, SMS received, scholarship, weekday) drive COUNTIFS and
SUMIFS formulas, so that the sheet works in any spreadsheet application. On the
\emph{Pivot Dashboard}, the same six criteria are slicers connected to every pivot table:
one click filters the KPI cards and the pivot charts at once. Each chart is anchored to a
fixed block of cells, so it keeps its place and size whatever the screen or the zoom.

\capture[0.74\linewidth]{dashboard.png}{Dashboard sheet: six KPI cards and eight charts,
all filters set to All.}

% =====================================================================
\section{Limitations}\label{sec:limits}

\textbf{Six weeks of data.} The file covers 29 April to 8 June 2016; no seasonality or
trend can be read.

\textbf{The reason for a no-show is not recorded,} nor the specialty or the time of day.

\textbf{The SMS flag says that a reminder was sent, not that it was read.} The SMS effect
is an observational estimate; a trial on the 1 to 3 day bookings would confirm it before
a full roll-out.

\textbf{Association, not causation.} The comparisons identify the groups at risk; they do
not predict an individual no-show.

\vspace{0.1cm}
\noindent\hfill\begin{minipage}{7cm}\begin{flushright}\singlespacing
\textbf{KOUAME Koffi Fid\`ele}\\
Data Analysis Internship\\
\href{mailto:koffifidelek59@gmail.com}{koffifidelek59@gmail.com}
\end{flushright}\end{minipage}

\end{document}
